% --- limitations ---
\paragraph{Class is confounded with acquisition session, and we tested whether it matters.}
The dataset was collected in class order: every whole-egg episode was recorded before every
cracked-egg episode (whole Dec~05 14:52--Dec~05 16:03, cracked Dec~05 16:22--Dec~06 15:38
UTC), with the class boundary falling at a session break. Any systematic change across that
boundary -- gel wear, membrane replacement, illumination from the sensor's internal LEDs,
mounting pressure, ambient temperature -- is correlated with the label and is in principle
available as a shortcut, and no partition of these episodes can separate the two. Rather
than leave that as a caveat we measured it.

Episode 11 is the one cracked egg recorded inside the whole-egg session, nineteen minutes after
the last whole-egg episode ends; every other cracked egg belongs to the following session.
Training on the whole-egg session plus the cracked session's other eggs and then testing on
that episode asks the question directly: a detector keyed to session appearance should fail
there, having seen only whole eggs from that session, while a detector keyed to crack
morphology should not. Against a control test set of cracked eggs drawn from the cracked
session with identical training data, the cross-session egg is detected equally well -- mean
AP$_{50}$ difference +0.011 over the four detectors, with recall between 0.778 and 1.000. The
shortcut is available but the models are not using it, so the cross-validated scores are not
explained by the confound.

That test rests on a single egg, nine frames, and that is the most the dataset permits: it
contains exactly one cracked episode outside the cracked session and no whole-egg episode
inside it, so no larger version of the probe and no reverse direction exist. A weaker
session effect would not be visible at this sample size. Re-acquiring a modest set of
episodes with the two classes interleaved within each session remains the way to settle it
-- now to confirm what one egg indicates, rather than to rule out a likely artefact.

\paragraph{Scale.} With 19 eggs the sample is small, and the grouped protocol makes that
explicit rather than hiding it: the spread across folds is wider than the differences
between architectures, which is why we quote a baseline range instead of a ranking. A single
split cannot expose that spread, and so invites exactly the ranking the data cannot support.
The absolute numbers here should be read as evidence about the sensing modality, not as a
benchmark result.
